% !PRINTED 470
\weberpage{470}

%<<CONT
oder genauer gesagt, je nach der Wahl des Vorzeichens in zwei Gleichungen, die außer dem Produkt
\[
\sqrt{2}\, f(\omega)^2
\]
nur rationale Zahlkoeffizienten enthalten. Letzteres ersieht man daraus, daß in der Gleichung (7) nur solche Produkte $u^h v^k$ vorkommen, in denen $h+k$ gerade ist (vgl. den Anfang von § 74).

Die Relation (14) fordert nun, daß $\omega$ einer Gleichung
\weq{16}{\frac{c + \partial\omega}{a} = \frac{\gamma(\omega+1) + \delta(\omega-1)}{\alpha(\omega+1) + \beta(\omega-1)}}
genüge, in der $\begin{pmatrix} \alpha, & \beta \\ \gamma, & \delta \end{pmatrix}$ eine zur ersten oder zweiten Klasse gehörige lineare Transformation ist. Damit aber eine der beiden Gleichungen (14) wirklich erfüllt sei, ist nach § 40, (12) noch erforderlich:
\weq{17}{(\alpha - \beta)(\alpha + \beta + \gamma - \delta) \equiv 0 \quad (\mathrm{mod}\ 8).}
Nun folgt aber, wenn $\omega$ Wurzel der quadratischen Gleichung (2) ist, aus (16) wie oben:
\weq{18}{
\begin{aligned}
\partial(\alpha + \beta) &= Ax, \\
\partial(\alpha - \beta) &= Bx + y, \\
c(\alpha - \beta) - a(\gamma - \delta) &= Cx, \\
c(\alpha + \beta) - a(\gamma + \delta) &= Bx - y, \\
2n &= y^2 + mx^2.
\end{aligned}
}
Ist nun, wie vorausgesetzt war,
\[
m \equiv 1 \quad (\mathrm{mod}\ 8),
\]
so setzen wir:
\[
n = \frac{m+1}{2} \quad \text{oder} \quad \frac{m+9}{2},
\]
je nachdem $m \equiv 1$ oder $\equiv 9\ (\mathrm{mod}\ 16)$ ist, so daß auch $n \equiv 1\ (\mathrm{mod}\ 8)$ wird, dann ist $x = 1$, $y = \pm 1$ oder $= \pm 3$ zu setzen, und wenn wir $A$ relativ prim zu $n$ voraussetzen, so wird $\partial = 1$, $a = n$, und aus (18) folgt: % [sic] printed second equality omits y
\[
\begin{aligned}
\alpha + \beta &= A, & n(\gamma - \delta) &= -C + c(B+y), \\
\alpha - \beta &= B + y, & n(\gamma + \delta) &= -B + y + cA.
\end{aligned}
\]
Hiernach ergibt die Bedingung (17), da $B+y$ ungerade, $c \equiv 0\ (\mathrm{mod}\ 16)$ ist:
\[
A \equiv C \quad (\mathrm{mod}\ 8),
\]
%CONT>>
